\begin{abstract}
As the deployment of specialized, fine-tuned deep learning models on resource-constrained edge devices becomes ubiquitous, post-training quantization (PTQ) to low-bit integers (e.g., INT8 or INT4) is essential to meet strict memory, storage, and thermal budgets.
Simultaneously, multi-task model merging is a vital paradigm for consolidating multiple single-task experts into a unified model without the high cost of multi-task training.
However, practitioners face a challenging conflict: merging experts in full precision and then quantizing the result introduces severe representation and alignment noise, while merging pre-quantized models violates linear mode connectivity.
To resolve this, test-time adaptation (TTA) approaches (e.g., AdaMerging, Q-Merge) optimize layer-wise coefficients, but they suffer from the \emph{Overfitting-Optimizer Paradox}---on tiny, unlabeled on-device calibration streams, unconstrained optimization over high-dimensional layer-wise spaces easily fits transductive noise, producing highly jagged coefficient schedules that fail to generalize.
We propose \textbf{Q-PolyMerge}, a hybrid-precision (weight-quantized, activation-float) framework that restricts the search space of merging coefficients to a low-dimensional continuous polynomial subspace of normalized layer depth.
By reducing the parameter space by over 78\% (e.g., from 56 to 12 parameters) and acting as a robust low-pass filter, Q-PolyMerge acts as a smooth continuous regularizer, effectively mitigating the risk of degenerate entropy collapse and overfitting.
Consequently, it stabilizes gradient-based optimization via the Straight-Through Estimator (STE) and enables zero-order derivative-free search (1+1 Evolution Strategy) to run efficiently with zero activation caching on edge microcontrollers.
Evaluating on a Vision Transformer (ViT-Tiny) across four benchmarks (MNIST, FashionMNIST, CIFAR-10, SVHN) under 8-bit and 4-bit PTQ, Q-PolyMerge significantly stabilizes adaptation. Under 4-bit PTQ, Q-PolyMerge (Adam STE) achieves an average accuracy of \textbf{48.87\%}, outperforming unconstrained Q-Merge by \textbf{+2.85\%} and naive post-merge quantization by \textbf{+5.95\%}. Under 8-bit PTQ, it achieves \textbf{59.76\%} average accuracy (nearly recovering the unquantized continuous PolyMerge ceiling of \textbf{61.00\%}) with significantly reduced variance, demonstrating smooth, physically stable coefficient schedules.
\end{abstract}
